\chapter{Orientation as Measurement}

The book has one thing left to settle, and it is the thing it has been postponing since
the pair was produced: the sign. The single invariant became a charge; the charge came in
a signed pair, an electron reading negative one and a positron reading positive one; and
the apparatus chose the electron's orientation and carried it without ever saying why that
orientation rather than its mirror. This final chapter takes up the sign, and it resolves
it in the book's own characteristic way --- not by deriving it, which cannot be done, but by
measuring it, which can. The sign convention is the book's last instance of its first
lesson: a measurement can anchor what it cannot derive.

The chapter opens by separating the two orientations cleanly. The electron and the positron
are not two charges; they are one residue read in two orientations, mirror images that
cancel. Reading the residue as measurement minus story gives negative one; reading it the
other way gives positive one; and their sum is zero. The sign is an orientation of the
reading, not a property of the residue's size, and the first section makes that distinction
exact.

The second section sharpens the distinction into an impossibility. The energy the apparatus
built carries magnitude and only magnitude. Its diagonal is nonnegative by construction ---
that was the whole content of positive-definiteness --- so a negative sign can never be a
diagonal energy reading. The sign does not live in the interior energy; it lives in the
cross term and in the orientation of the instrument. The interior is sign-blind, and the
section proves that blindness is structural: from the energy alone, the apparatus can read
how much charge is present but never which orientation it carries. The sign is not waiting
in the interior to be derived. It is not there.

The third section says what the apparatus can do instead, and it is the oldest move in the
book: count. If the sign cannot be derived from the interior, it can be anchored from the
outside, by tallying which orientation occurs and asking whether one clears the other by a
margin. The apparatus reuses the separation logic it has used since it first grounded
comparison --- a candidate becomes anchored when its measured separation clears a noise
floor --- and points it at orientation. One orientation dominates when its count exceeds the
other's by more than the floor. Dominance is the measurement that anchors what derivation
cannot reach.

The fourth section proves the anchor is well-defined. Two orientations cannot both dominate
the same tally; dominance, once it holds, holds for exactly one orientation. So a convention
anchor --- a tally together with its dominant orientation and the proof that it dominates ---
is data carrying its own certificate, not a choice dressed up as a fact. Anchors built on
the same tally agree on the dominant orientation, and a concrete sample shows the machinery
turning: a count of seven against two, past a floor of two, anchors the electron, and the
anchored reading is negative one.

The fifth section states the chapter's --- and the book's --- closing thesis. Convention is
gauge; tally is measurement. The apparatus does not derive the electron's sign, and it does
not pretend to; it anchors the sign by counting, taking ``one orientation occurs more'' as
empirical input and reading the dominant orientation off the tally. This is the book's whole
method in one last instance: the apparatus sits in a convention it cannot derive, names it
as the convention it is, and anchors it by a measurement rather than smuggling it in as a
hidden assumption. The book has carried an exposed decision edge from its first
pages --- the single classical edge the apparatus sits in rather than computes ---
and this chapter names and anchors that edge as sign.

By the end the book is closed and honest about its own edge. The single invariant has become
a charge, the charge a signed pair, and the sign an anchored convention --- anchored by a
tally, not derived from the interior, with the deeper question of why one orientation should
dominate left, explicitly, to a physics beyond the finite apparatus. The book set out to
measure a single invariant. It measured the invariant, produced the electron, generated the
gauge, read the radiation, ordered its set theory, named the continuum door, and now anchors
the one sign it could never derive --- closing, as it began, by sitting deliberately in the
one decision it cannot compute.

\section{Orientation and Opposite Orientation}

Chapter 16 left the book at a threshold. The finite apparatus had reached the
completion door honestly: it had named the conditions under which the boundary
reading would pass into an analytic setting, and it had refused to pretend that
the closing had already occurred. That restraint leaves one last finite question
on the table. The apparatus has carried a signed charge since the pair was
produced, but it has not yet said what kind of fact the sign is. Is the negative
orientation an interior output, or is it an orientation the apparatus must
anchor by measurement?

The first step is to get the two orientations exactly right. The temptation is
to treat the electron and the positron as two unrelated charges, two separate
objects that happen to balance. That is not the structure built here. The pair
comes from one residue read in two directions. The same unit difference can be
read from the measuring side against the story, or from the story side against
the measurement. Reversing the direction reverses the sign. It does not create
a second magnitude.

Recall the production rule. The pair kick gave a coupled measurement and a
linear story, and the residue was their difference. Read as measurement minus
story, the residue is negative one. Read as story minus measurement, the same
residue is positive one. The two readings are not two independent quantities;
they are one quantity read with two orientations of subtraction:
\[
  \begin{aligned}
    \text{electron} &= \text{measurement} - \text{story} = -1,\\
    \text{positron} &= \text{story} - \text{measurement} = +1.
  \end{aligned}
\]
The signs are opposite because the readings face opposite ways. The electron
reading is not a different residue from the positron reading. It is the same
residue seen with the instrument pointed in the other direction. Neither reading
is more arithmetically real than the other; each is what the same finite record
returns when the subtraction is oriented from one side.

This matters because a sign can look like extra substance when the orientation
has been forgotten. If one begins with ``electron'' and ``positron'' as names
already carrying physical authority, it is easy to imagine two independently
stored charges. But the finite record is simpler and stricter. It stores the
unit residue and the direction in which that residue is read. The name
``electron'' marks the measurement-minus-story orientation. The name
``positron'' marks the opposite orientation. The charge sign is the direction of
the reading.

Because they are one residue oriented two ways, the two readings cancel. Add
the orientations and the residue subtracts from itself:
\[
  \text{electron} + \text{positron} = (-1) + (+1) = 0.
\]
The cancellation is not a coincidence to be explained by a balance of two
populations. It is the algebraic fact that a quantity minus itself is zero,
read once forward and once backward. The pair, summed, is silent because the
pair is one residue in two orientations and the two orientations are negatives
of each other. The zero is not an erasure of the charge; it is the signature of
opposite readings being added together.

This is why the apparatus speaks of orientation rather than of two charges. An
orientation is a choice of which way to read a difference, and the residue
admits two. The electron orientation reads measurement against story. The
positron orientation reads story against measurement. Switching orientation
flips the sign without touching the unit itself. The charge has one magnitude,
one unit, and two orientations, exactly as a difference has one size and two
signs.

The section is careful to keep orientation distinct from magnitude even in
language, because the rest of the chapter turns on that distinction. Magnitude
answers the question how much. Orientation answers the question which way. The
unit residue answers the first question: the charge is one unit in size. The
choice between the two readings answers the second: electron-side or
positron-side, measurement-minus-story or story-minus-measurement. A signed
reading reports both, but the two parts do not have the same status.

The magnitude belongs to the finite residue that has already been produced.
The apparatus can display it as a unit difference. The orientation belongs to
the way the record is read. The apparatus can display both directions, but
displaying both is not the same as selecting one as the convention it will
carry. The mirror reading is always available. If the electron orientation is
called negative, the opposite reading is positive; if the labels were reversed,
the algebraic opposition would remain. What would change is the convention of
which side the apparatus treats as electron-side.

That last sentence is the hinge of the chapter. The apparatus does not need to
invent a second charge in order to have a positron reading. It only needs the
opposite orientation of the same unit residue. Nor does it need to pretend that
the negative sign is stored in the interior energy. At this point the section
has done only the preliminary work: it has separated the unit from the direction
in which the unit is read. Magnitude and orientation have been named as
different aspects of one signed reading.

The later sections will ask what can anchor the orientation convention. They
will not ask the finite interior to do a job it cannot do, and they will not
appeal to a physical asymmetry theorem outside the apparatus. They will ask a
measurement question: when the apparatus is given occurrences of the two
orientations, can one orientation dominate past a noise floor, and can that
dominance serve as the anchor for the convention? This section prepares that
question by making the two candidates precise. There are two orientations of
one residue, not two unrelated charges; their magnitudes agree, their signs
oppose, and their sum is zero.

So the chapter begins with a modest but necessary correction. The electron and
the positron are not a particle and its rival in this finite account. They are
one residue read from its two sides. The sign is not the size of the residue.
It is the orientation of the reading. The book's final task is to say how an
apparatus can anchor which side it calls negative while remaining honest that
the opposite side is present, coherent, and algebraically equal in magnitude.

\section{Magnitude versus Sign}

The previous section separated magnitude from orientation. It showed that the
electron and the positron are one unit residue read from opposite sides, not two
unrelated charges. This section tightens that separation into a structural
claim: the apparatus can read the size of the charge from its interior energy,
but it cannot read the sign there. The sign belongs to the oriented reading of
the residue. The size belongs to the nonnegative interior account.

That distinction can feel too small, because in ordinary notation a signed
number presents magnitude and sign at once. The symbol \(-1\) seems to arrive as
one item: a unit together with a direction. The finite apparatus has to be more
careful. It must ask which part of the signed reading comes from the energy it
has built and which part comes from the way an instrument reads a difference.
The answer is not optional. The diagonal interior energy has the right shape to
report magnitude and the wrong shape to report sign.

Recall what positive-definiteness says. The energy is the diagonal of the
pairing, and that diagonal is nonnegative for every activity, vanishing only at
the trivial one:
\[
  E^2(x) = a(x,x) \ge 0,
  \qquad
  E^2(x) = 0 \Leftrightarrow x = 0.
\]
Earlier chapters needed this property because a detector that could assign a
negative amount of its own activity would fail to distinguish presence from
absence. A genuine diagonal energy must make activity visible as size. If the
activity is absent, the size is zero. If the activity is present, the size is
positive. The diagonal therefore does exactly what the apparatus needs from an
interior detector: it turns a finite activity into a quantity that can be
compared, bounded, and carried forward without ambiguity about whether anything
has occurred.

Read the same property now for what it forbids. A nonnegative quantity can tell
the apparatus how much. It cannot tell the apparatus which way. It has no place
for a minus sign as an output value. If the charge has unit size, the diagonal
energy can display that unit size. If the charge is read in the electron
orientation, the diagonal still displays the same unit size. If the charge is
read in the positron orientation, the diagonal again displays the same unit
size. Nothing in the diagonal changes from \(+1\) to \(-1\), because the
diagonal never returns a negative value in the first place.

This gives the apparatus a constructive magnitude statement rather than a mere
prohibition. The apparatus can read the unit size of the residue. It can say
that the produced charge is not zero. It can compare that size with other finite
sizes and hold it as a stable part of the record. Magnitude is therefore not a
failure of the apparatus. The finite interior succeeds at reading magnitude
precisely because the diagonal energy is nonnegative and definite.

But success at magnitude is also blindness to orientation. The same
positive-definite form that makes the detector reliable removes sign from the
detector's range. The interior does not hide a choice between electron and
positron behind the unit reading. It does not contain a concealed negative value
waiting for a sharper calculation. It gives the same nonnegative answer for the
two opposite orientations because the two orientations have the same size. The
apparatus can extract \(|\text{charge}| = 1\) from the interior account; it
cannot extract which side of the residue receives the negative name.

The difference lives where the previous section placed it: in the oriented
residue. The residue comes from a mixed comparison, not from an activity paired
with itself. A diagonal term \(a(x,x)\) asks one activity for its size. A mixed
term asks how two channels face one another. Once two channels enter the reading,
orientation matters. Measurement-minus-story and story-minus-measurement are
not the same directed reading, even though they have the same unit size. They
are negatives of one another because the instrument has reversed the order in
which it reads the difference.

That is the role of the cross term. A cross term can carry a sign because it is
not the diagonal energy of a single activity with itself. It is a coupling
between channels, and a coupling can change sign when the order or orientation
of the reading changes. The apparatus does not ask the cross term to behave like
a nonnegative storehouse. It asks the cross term to record how the two sides
meet. In that meeting, the same unit residue can face one way or the other.
The sign is the orientation of that meeting as read by the instrument.

So the two responsibilities divide cleanly:
\[
  \begin{aligned}
    \text{diagonal energy} &\longmapsto \text{nonnegative magnitude},\\
    \text{oriented cross reading} &\longmapsto \text{signed residue}.
  \end{aligned}
\]
The first line reports the size of the activity the apparatus has produced. The
second line reports the direction in which the produced residue is read. Neither
line replaces the other. If one tries to make the diagonal line carry sign, the
positivity of the diagonal blocks the attempt. If one tries to make the cross
reading carry magnitude without orientation, the actual electron/positron pair
has been flattened into a signless size and the distinction from the previous
section disappears.

This is the precise sense in which the interior is sign-blind. Hand the
interior an activity and it returns a magnitude with perfect reliability and no
orientation whatsoever. The interior cannot tell the electron reading from the
positron reading, because the only thing it contributes at this level is a
nonnegative energy, and the two readings have the same nonnegative energy. The
difference between them is not an interior energy fact. It is an orientation of
the residue, and orientation is not something a nonnegative reading can hold:
\[
  \text{interior reads } |\text{charge}| = 1,
  \qquad
  \text{interior cannot read } \mathrm{sign}.
\]
The blindness is structural, not a limitation of effort. No refinement of the
same diagonal energy can make it produce a negative sign while it remains the
positive-definite diagonal energy. Refinement can sharpen a size, reduce an
error, or improve a comparison among nonnegative values. It cannot turn the
range of the diagonal into two oriented signs without changing what the diagonal
is. The apparatus should not ask a magnitude detector to perform the work of an
orientation reader.

This also protects the sign from a subtler mistake. One might say that the
interior energy chooses the sign indirectly: perhaps the energy itself stays
nonnegative, but some hidden preference in that energy tells the apparatus
which orientation to call electron. The finite account refuses that move. A
nonnegative value does not secretly contain one of its two oriented square
roots as a favored convention. The value \(1\) supports the two signed readings
\(+1\) and \(-1\), but the value \(1\) alone does not tell the apparatus which
reading should carry which name. The diagonal can certify the unit; it cannot
name the orientation.

Nor does the apparatus smuggle the sign into the word ``charge.'' If the word
already meant ``negative electron orientation,'' the chapter would have solved
the problem by vocabulary. Instead, the apparatus keeps the word under
discipline. A signed charge report contains a magnitude and an orientation. The
magnitude comes from the finite residue's size. The orientation comes from the
directed reading of that residue. Until an orientation has been anchored, the
apparatus may display both readings and may call them opposite, but it has not
explained why one convention should govern the final naming.

The point can be seen through cancellation. The two oriented readings add to
zero:
\[
  (-1) + (+1) = 0.
\]
The diagonal energy does not see this cancellation as a negative quantity
destroying a positive one inside the energy account. It sees equal size on both
sides. Cancellation belongs to the oriented signed readings, where one reading
faces the other. The interior energy keeps the size of each reading visible,
but the algebraic opposition occurs in the oriented residue ledger. That is why
the pair can cancel as signs while sharing one magnitude.

This separation keeps the apparatus honest about what it has and has not
measured. It has measured a unit. It has measured a finite difference whose two
directions oppose. It has not measured, from the diagonal interior alone, a
reason to privilege one direction as the electron convention. The diagonal
energy supplies no such reason. It neither prefers the measurement-minus-story
orientation nor prefers the story-minus-measurement orientation. It simply
reports the size common to both.

The chapter therefore cannot close by saying that the sign has been computed
from the interior. The interior computation has reached its natural boundary.
It gives a magnitude certificate: activity is present, the residue has unit
size, and the two oriented readings share that size. It does not give an
orientation certificate. A certificate for orientation must include the
directional fact itself, not merely the nonnegative size that both directions
share.

This is why the sign belongs to an instrument reading. The instrument reads a
residue in an order. It can face the record from measurement toward story or
from story toward measurement. Once it faces the record, the sign appears as
the direction of that reading. The instrument can reverse, and when it reverses
the sign reverses. The finite apparatus can keep both possibilities in view
without contradiction because the underlying magnitude remains the same.

The result is not skepticism about sign. It is a more precise account of sign.
The sign is real as an oriented reading; it is not a diagonal energy value. The
electron reading is not less real because the energy did not output a negative
number. It is the measurement-minus-story orientation of the produced residue.
The positron reading is not less real because it mirrors the electron. It is
the opposite orientation of the same produced residue. What the interior cannot
do is decide, from its nonnegative magnitude alone, which orientation the
apparatus should take as its governing convention.

That undecided edge is exactly where the next section begins. Once magnitude
and sign have been separated, the remaining question is no longer ``Where in
the interior does the minus sign hide?'' The remaining question is ``How can an
apparatus anchor an orientation convention without pretending that the
convention came from diagonal energy?'' The book's answer will stay finite. It
will not appeal to an outside theorem selecting the electron orientation, and
it will not turn the nonnegative unit into a secret sign. It will ask the kind
of question a measuring apparatus can answer: given occurrences of the two
orientations, can one orientation dominate the other past a noise floor?

The present section stops just before that tally. It has done the necessary
negative and positive work. Negatively, it has ruled out the interior energy as
a source of sign. Positively, it has preserved the interior energy as the
source of magnitude. The apparatus reads how much through the nonnegative
diagonal, reads which way through the oriented residue, and must anchor its
final convention by a measurement suited to orientation rather than by a
magnitude detector asked to speak beyond its range.

\section{Dominance past a Noise Floor}

The previous section left the apparatus with a clean division of labor. The
interior energy can read the unit size of the residue, but it cannot decide the
orientation convention. The sign is not a diagonal energy output. It is the
direction of an oriented reading. If the apparatus wants to anchor that
direction, it needs a measurement suited to orientation rather than a
magnitude detector asked to speak outside its range.

The measurement suited to orientation is a tally. The apparatus counts
occurrences of the two possible readings: the measurement-side orientation and
the story-side orientation. For the purposes of the chapter, it may name these
the electron orientation and the positron orientation. The names do not explain
why either occurs. They mark the two directions whose frequencies the apparatus
can compare.

A tally is not a derivation of sign. It does not say why one orientation
appears more often than the other. It does not reach into the interior and find
a missing minus sign. It does something more modest and more useful: it gives
the apparatus a finite record of occurrence. Where the diagonal energy could
only report a nonnegative unit, the tally reports how many times each oriented
reading has appeared. The apparatus can hold those counts, compare them, and
refuse to say more than the comparison warrants.

That last refusal matters. Raw difference is not yet dominance. If one
orientation appears ten times and the other appears nine times, the apparatus
has a difference, but it may not have a trustworthy anchor. A small separation
can belong to noise, sampling accident, or unresolved fluctuation. The
apparatus would smuggle convention in by fiat if it treated every nonzero
difference as decisive. It needs the same discipline it has used throughout the
book: a measured separation must clear a floor before the apparatus treats it
as an anchor.

The noise floor is not a physical theorem selecting the electron. It is a rule
of measurement discipline. Below the floor, the apparatus declines the call.
It does not pretend that close counts have settled the convention. It says:
these two orientations remain too close for this tally to anchor either one.
That refusal is not failure. It is how the apparatus preserves honesty when
measurement has not separated the candidates strongly enough.

Above the floor, the situation changes. If one orientation exceeds the other by
more than the floor, the apparatus has a measured dominance relation. The
dominant orientation is not magically produced by the count; it is identified
by the count. The tally says which orientation occurs more often by a margin
the apparatus has agreed to trust. That margin turns a raw difference into a
candidate convention anchor.

The rule can be written without metaphysical decoration. Let \(o\) be one of
the two orientations. Let the other orientation be its mirror. Then \(o\)
dominates exactly when the count of \(o\) exceeds the count of the other
orientation by more than the noise floor:
\[
  \text{$o$ dominates}
  \quad\Longleftrightarrow\quad
  \mathrm{count}(\text{other orientation}) + \mathrm{noise\ floor} < \mathrm{count}(o).
\]
This formula gives the apparatus three possible outcomes, not two. The electron
orientation may dominate. The positron orientation may dominate. Or neither may
dominate because the counts lie within the floor. The third outcome is
essential. Without it, the apparatus would always choose a convention, even
when the measurement gave it no right to choose. The noise floor creates a
place for principled silence.

Principled silence keeps the sign from becoming a hidden derivation. Suppose
the counts fall below the floor. The apparatus still has both oriented
readings. It still has their common magnitude. It still has a finite tally. But
it does not have a convention anchor. It must leave the orientation convention
unsettled for that tally. The apparatus has measured and found no dominance.
That result carries information: the finite record does not support a
selection.

Suppose instead that one orientation clears the floor. Then the apparatus may
use the dominant orientation as the measured candidate for the convention. It
does not thereby explain why the world presented that orientation more often.
The count supplies no origin story. It supplies a finite anchor: this
orientation occurred more than its mirror by a margin the apparatus recognizes
as decisive. That is enough for a convention anchor, and it is exactly the kind
of enough the book has been defending from the first mark onward.

This is gauge fixing by measurement. The interior left the orientation
undetermined, not because the sign is unreal, but because diagonal magnitude
cannot select a direction. The tally fixes that freedom only where occurrence
has separated one orientation from the other beyond the floor. The convention
is gauge; the tally is measurement. The apparatus does not invent the sign, and
it does not compute it from energy. It reads dominance where dominance appears
and refuses the convention where dominance does not appear.

The language of dominance must remain narrow. Dominance means one count exceeds
the other count by more than the floor. It does not mean that the dominant
orientation has metaphysical priority over its mirror. It does not mean that
the other orientation is false. It does not mean that an external asymmetry
theorem has entered through the back door. It means only that this finite tally
has separated the two orientations enough for the apparatus to anchor one as
the convention it will carry.

This narrowness also prepares the next question. If a tally can anchor a
convention only when one orientation dominates, can both orientations ever
dominate the same tally? If both could dominate, the anchor would collapse into
ambiguity. The next section answers that uniqueness question. The present
section does not prove uniqueness yet. It sets the condition whose uniqueness
will matter: dominance past a floor, measured on a finite tally of the two
oriented readings.

The result is the chapter's turn from impossibility to measurement. S02 showed
that interior magnitude cannot carry sign. This section shows what the
apparatus can do instead. It counts occurrences, compares the counts, applies a
noise floor, refuses close calls, and treats only floor-clearing dominance as a
candidate convention anchor. The sign does not come from the interior. It is
anchored, when it is anchored, by a count that has earned the right to speak.

\section[The Convention Anchor]{The Convention Anchor and the Unique\\ Dominant Orientation}

A measurement that anchors a convention must anchor it definitely, or it has
not anchored a convention at all. The previous section gave the apparatus a
tally and a noise floor. It also gave the rule: one orientation dominates when
its count exceeds the other's by more than the floor. This section proves that
the rule cannot point in both directions at once. For a fixed tally and a fixed
floor, at most one orientation can dominate. That uniqueness is what turns a
dominance reading into a convention anchor rather than a preference with a
number attached.

Keep the frame finite. The apparatus has two counts, one for each orientation.
It also has a nonnegative floor, the margin below which it refuses to call a
separation decisive. Nothing else enters the proof. No hidden physical
asymmetry selects the outcome, and no interior energy supplies a missing sign.
The apparatus asks only whether the two finite inequalities that would make the
two orientations dominate can hold together.

Suppose, for contradiction, that they could. Suppose the electron orientation
cleared the positron count by more than the floor, and suppose the positron
orientation also cleared the electron count by more than the same floor:
\[
\begin{aligned}
  \mathrm{count}(\text{positron}) + \text{floor}
    &< \mathrm{count}(\text{electron}),\\
  \mathrm{count}(\text{electron}) + \text{floor}
    &< \mathrm{count}(\text{positron}).
\end{aligned}
\]
Add the inequalities. The electron count and the positron count appear on both
sides, once as the larger side and once as the smaller side. After cancellation,
the two assumptions force
\[
  2\cdot \text{floor} < 0.
\]
But the floor is nonnegative. Twice a nonnegative floor cannot fall below zero.
So the two dominance claims cannot both hold. A fixed tally at a fixed
nonnegative floor admits at most one dominant orientation.

This is the whole uniqueness proof. It is deliberately elementary because the
anchor it protects is elementary. The apparatus has not discovered a deep law
about why one orientation appears more often. It has shown that its own
dominance rule is single-valued: when the tally clears the floor for one
orientation, the mirror orientation cannot also clear the floor. The rule may
produce one dominant orientation or no dominant orientation. It never produces
two.

That last distinction matters. Uniqueness does not say that every tally yields
an anchor. A tie yields none. Counts separated by less than or equal to the
floor yield none. A tally that does not clear the floor in either direction
leaves the convention unanchored for that record. The apparatus must keep the
third outcome alive: no anchor. Only when one orientation clears the floor does
the apparatus get data strong enough to carry a convention.

So a convention anchor is not merely a preferred orientation. It is a finite
package:
\[
  \text{anchor}
  =
  (\text{tally},\ \text{floor},\ \text{dominant orientation},\
    \text{dominance certificate}).
\]
The tally records the two counts. The floor records the margin of refusal. The
dominant orientation names the side that cleared the floor, when such a side
exists. The certificate records the inequality showing that the named side
really cleared the mirror side by more than the floor. If no such inequality
holds, the package does not exist.

This definition turns the convention into data with a boundary. It says exactly
what the apparatus may carry and exactly when it must stop. If the data include
no floor-clearing orientation, the apparatus cannot manufacture one by taste.
If the data do include such an orientation, the apparatus can carry the
orientation together with the tally and margin that witness it. The convention
anchor is therefore a measured record, not an interior verdict and not a
reader's preference.

The certificate also keeps the anchor auditable. A reader can inspect the tally
and floor and check the inequality directly. The reader need not trust the
apparatus's temperament, its naming habit, or its desire for the electron
orientation to win. The data say whether the proposed orientation dominates.
If the proposed side fails the inequality, the anchor fails. If the proposed
side passes it, the mirror side cannot pass as well.

The margin is part of the certificate, not decoration. If the electron count is
seven, the positron count is two, and the floor is two, the margin is five and
the excess past the floor is three. The apparatus can point to that excess as
the finite reason the call is allowed. If the excess vanishes, or if it merely
equals the floor, the apparatus withholds the anchor. This keeps the convention
from sliding into a bare majority rule. The apparatus does not ask which count
is larger; it asks which count is larger by enough.

That ``by enough'' belongs to the record because the floor belongs to the
record. A tally without its floor is incomplete for convention anchoring. The
same pair of counts may anchor under one floor and fail under a stricter floor.
For example, seven against two clears a floor of two, but it would not clear a
floor of five. The apparatus therefore never reports just ``seven against
two.'' It reports the tally together with the floor under which the comparison
has meaning. The convention anchor carries both.

This is why the anchor differs from an ordinary label. A label can be copied
without its evidence. An anchor cannot. To carry the anchor is to carry the
finite conditions under which the orientation dominates. If a later reader
changes the floor, the reader has changed the measurement question and must ask
again whether dominance holds. If the reader keeps the same tally and floor,
the uniqueness proof keeps the answer fixed.

This is the sense in which the convention becomes independent of the reader.
Give the same tally and the same floor to two readers. If neither orientation
dominates, both readers must withhold an anchor. If one orientation dominates,
both readers must identify the same orientation, because the two-inequality
argument rules out any second dominant side. The reader may choose different
words, but the finite relation in the tally stays fixed.

Independence here does not mean independence from measurement. It means
independence within measurement: once the finite record and the floor are fixed,
the result no longer depends on the hand that reads them. The apparatus does
not float above the tally and pronounce a sign. It submits its convention to a
finite check. The check either produces a unique dominant orientation or
produces no anchor.

The no-anchor case deserves the same respect as the positive case. If neither
orientation dominates, the apparatus has not failed to be decisive; it has
successfully obeyed its own measurement rule. It keeps both orientations
available and refuses to promote either one into the governing convention for
that tally. That refusal protects the later convention from overreach. A
measured convention must emerge only where the measurement has enough
separation to support it.

A concrete sample shows the machinery turning. Tally nine events: seven in the
electron orientation and two in the positron orientation. Let the noise floor
be two. The electron orientation dominates because the positron count plus the
floor remains below the electron count:
\[
  \mathrm{count}(\text{positron}) + \text{floor} = 2 + 2 = 4 \;<\; 7 = \mathrm{count}(\text{electron}).
\]
The anchor, in this example, names the electron orientation. The associated
signed reading is the electron reading, negative one. The example does not say
why seven events occurred in that orientation and two in the other. It does not
claim a physical law has chosen the electron. It only illustrates the
dominance rule: with these counts and this floor, the electron orientation
clears the margin and the positron orientation cannot also clear it.

Change the sample and the apparatus behaves differently. If the tally were
five and four with the same floor of two, no anchor would appear. The larger
count would not exceed the smaller count by more than the floor. The apparatus
would have a finite record, but not a convention anchor. If the tally were two
and seven, the mirror orientation would dominate instead. The rule does not
prefer the electron in advance. It follows the measured separation.

This keeps the convention honest. The apparatus may carry a sign convention
only when the tally provides the appropriate certificate. It may not appeal to
the positive-definite interior, because that interior reads magnitude and stays
silent about orientation. It may not appeal to the mere fact that one count is
larger, because below the floor the apparatus refuses close calls. It may not
appeal to taste, because the certificate must be inspectable in the finite
record.

Notice also that the anchor never erases the mirror orientation. Even when the
electron orientation dominates in a tally, the positron orientation remains the
opposite reading of the same unit residue. The anchor does not say that the
mirror has become incoherent or unreal. It says that the finite tally has given
the apparatus a certified reason to carry one side as the convention. The
opposite side remains available as the opposite orientation; it simply does not
win the dominance comparison under that record and floor.

This distinction matters for the book's final honesty. A convention anchor is
stronger than arbitrary naming, because it comes with a tally, floor, dominant
orientation, and certificate. It is weaker than an explanation of why the world
offered that tally, because it does not supply such an explanation. The anchor
occupies the middle position the apparatus has earned: finite, inspectable,
unique when it exists, deliberately recorded, and silent about origins beyond
the measured record.

The section therefore proves exactly what S03 left open. S03 defined
dominance. S04 shows that dominance, when it exists for a fixed tally and
floor, is unique. From that uniqueness the apparatus can form a convention
anchor as data: not just ``electron,'' but this tally, this floor, this
orientation, and this margin. The anchor comes with its own reason for being
carried and its own boundary against overuse.

This prepares the closing section. The book still must say what kind of
success this is. A unique dominance anchor lets the apparatus carry a
convention without mere preference, but it does not turn convention into an
interior theorem. It anchors the sign by measurement rather than pretending the
sign came from magnitude. The convention is now finite data with a unique
orientation when the tally clears the floor; the final section will name that
achievement without asking it to be more than it is.

\section{Anchoring a Convention without Deriving It}

The book closes on a single sentence, and the chapter has earned the right to say it
plainly: convention is gauge, and tally is measurement. The sign the apparatus carried from
the moment it produced the pair is a convention --- a gauge freedom the interior left open ---
and the apparatus fixes it not by deriving it but by measuring which orientation the world
presents more often. This last section says exactly what that does and does not accomplish,
because the honesty of the whole book lives in the distinction.

What it accomplishes is an anchor. The apparatus has a tally, a floor, a
dominant orientation when one exists, and a certificate saying by how much the
dominant side clears the comparison. Those four pieces are the anchor. They do
not float as a slogan above the record. They are the record in the form needed
to carry a convention: this finite count, this refusal of close calls, this
orientation, and this margin. When the record has that shape, the apparatus can
say which side of the residue it will carry as negative. It can do so without
pretending that preference alone settled the sign, because the convention now
has a measured support.

The anchor therefore fixes the convention conditionally. If the tally says that
the electron orientation clears the positron orientation past the floor, then
the apparatus carries electron-orientation-as-negative. If the opposite tally
clears the same floor, then the opposite orientation would be carried instead.
If neither count clears the floor, the apparatus produces no convention anchor
from that record. This last case is not a failure of nerve. It is part of the
definition. A finite apparatus that cannot distinguish a close comparison past
its own floor has no right to smuggle a decision into the gap.

The achievement is reader-independent in exactly this finite sense. Any reader
given the same tally and the same floor reads the same result, because S04 has
already shown that two opposite dominant orientations cannot both be certified
for that fixed record. The reader need not share the apparatus's taste, history,
or preferred name for the sign. The reader only checks the counts, the floor,
and the margin. If dominance exists, the same orientation is singled out. If
dominance does not exist, the same absence of an anchor is visible. That is the
kind of objectivity the chapter can honestly claim.

What the anchor does not accomplish is just as important. It does not explain
why the world supplied that tally. It does not say why one orientation appears
more often in the measured record. It does not turn the sign into an interior
property of the positive magnitude, and it does not replace any physics that
might later account for an asymmetry in occurrence. The apparatus receives the
finite count as input. It can check whether the count clears the floor; it can
use the cleared count to anchor a convention; it cannot manufacture the fact
that the events leaned one way rather than the other.

This boundary keeps the chapter from overclaiming at the moment when
overclaiming would be easiest. The apparatus has done real work. It separated
orientation from magnitude, showed why the interior cannot see the sign, gave a
finite rule for dominance, and proved that dominance is unique when it exists.
Those steps are enough to anchor a convention by measurement. They are not
enough to say that the convention was forced by the interior. The anchor is
stronger than arbitrary naming and weaker than an origin story. It is exactly
where finite measurement can stand.

Put differently, the chapter has changed the kind of ignorance the apparatus
has about sign. At the start, the sign could look like an arbitrary leftover: a
name attached to one side of a symmetric pair. After the tally rule, the
apparatus is no longer merely naming in the dark. It has a finite test for
whether the world has presented one orientation more strongly than the other,
and it has a finite rule for refusing to act when the presentation is too close.
The ignorance that remains is narrower and more honest. The apparatus still
does not know why the events have the distribution they have. It does know what
to do with the distribution once it is on the table.

This matters because a convention can be empty if it has no check on it. A
reader could always decide to call one side negative and the other positive,
but such a decision would not yet be a measured convention. It would be a label
with no receipt. The anchor supplies the receipt. It says: here is the tally;
here is the floor; here is the margin; here is the orientation that survives
the comparison. The convention is still a convention because the interior did
not force the sign. It is measured because the apparatus can point to the finite
record that licenses carrying this convention rather than leaving the sign
unseated.

Nor is the anchor a majority vote in disguise. The floor is essential. Without
it, a one-count difference could settle the sign even when the apparatus itself
has declared that such a difference should not be trusted. The floor records
the apparatus's discipline about noise and close calls. The margin records that
the discipline has been met. Together they prevent the tally from becoming a
decorated preference. The apparatus carries the sign only when the count has
passed the standard by which the apparatus agreed to be judged.

The mirror orientation remains visible throughout this move. When the electron
orientation is carried as negative, the positron orientation has not vanished;
it is the opposite side of the same signed residue. The anchor does not destroy
the symmetry of naming, and it does not say that the losing side is unreal. It
says only that, for this tally and this floor, one side is the measured
dominant side. That distinction lets the apparatus carry a definite convention
without turning the opposite convention into nonsense.

The same restraint governs repeated use of the anchor. A later record may have
to be checked on its own terms: its own tally, its own floor, its own margin,
and its own finite audit.
The apparatus may carry an accepted convention forward as part of its working
language, but the measurement claim always points back to the record that
licensed it. That keeps the convention from hardening into an unexplained law.
The sign is carried because a finite count supported it, not because every
future comparison has already been settled in advance. The anchor remains a
measured permission, not a permanent exemption from later finite measurement
review by later readers.

The decision edge that has accompanied the book from its first mark can now be
named without being inflated. Earlier chapters needed a place where the
apparatus sat in a choice rather than computing one away. This chapter identifies
the relevant instance: the orientation of the signed residue, the side the
apparatus will call negative. That does not mean the earlier needle secretly was
charge sign all along. It means the sign is the final, sharpest case of the same
structure: a finite construction can expose the place where a convention must
be seated, but it must not pretend that seating the convention is the same as
deducing it from what came before.

The one-place shape is therefore simple. There is one sign convention to carry,
and there is one measured anchor that may license carrying it. The apparatus
does not install two signs, one official and one hidden. It does not carry an
unmeasured sign under the name of measurement. It waits for the tally, applies
the floor, checks the margin, and seats the convention only when the record
licenses that seating. One convention, one anchor, one explicit boundary around
what has and has not been earned.

This is why the chapter's closing sentence matters. Convention is gauge: the
interior construction leaves the sign open as a choice of orientation. Tally is
measurement: the finite record can anchor that choice when one orientation
dominates past the floor. Together the two clauses say neither too little nor
too much. They do not reduce the sign to whim, because the tally may certify a
dominant orientation. They do not promote the tally into a theory of why the
dominance occurred, because the count remains empirical input.

There is a useful modesty in that sentence. It does not ask the tally to become
a metaphysics of occurrence. It asks the tally to do what tallies can do:
record how often each orientation appeared to the apparatus, relative to a
specified standard for calling the difference decisive. It does not ask gauge
language to hide the sign as if naming were enough. It asks convention language
to mark the place where the interior left a choice open. The two halves need
one another. Convention without tally would be bare naming; tally without
convention would be a count with no sign to carry. Together they give the
finite apparatus a disciplined way to move from observed occurrence to carried
orientation.

The section also closes the numbered argument of the chapter. S01 introduced
opposite orientations. S02 separated sign from magnitude. S03 made occurrence a
finite tally question. S04 proved uniqueness of dominance for a fixed tally and
floor. S05 now names the result: the apparatus can anchor a convention by
measurement without claiming to derive the sign. That is the whole logical
movement. The bridge can now look back over the book's carried needle, but this
numbered section has finished the chapter's argument.

The final metaphysical anchor is the book's plainest one. Measurement does not
erase every convention. Sometimes it records enough to carry a convention
honestly. The apparatus constructs what can be constructed, names what remains open, and
then asks the finite record whether that openness has a certified orientation.
When the answer clears the floor, the sign is anchored. When it does not, the
apparatus stays silent. Either way, the world is not explained away. It is
counted, and the convention is carried only as far as that count allows.

\section*{Bridge: The Needle, Named at Last}

This is the last bridge, and it has no next chapter to force. Its work is
different from the earlier bridges. They pressed the book forward from one
unfinished question into another. This one gathers the chapter, gathers the
needle the book has learned to recognize, and lets the numbered argument close.
The reader is not being pushed into a new construction. The reader is being
asked to see where the construction has finally stopped.

Chapter 17 began with two opposite orientations of one residue. The electron
and positron readings were not treated as unrelated things, nor as two
independent units competing for the same place. They were two ways of reading a
single signed residue from opposite sides. That first step mattered because a
sign question cannot be handled honestly until the two sides of the sign have
both been kept in view. The apparatus had to preserve the mirror before it
could ask whether one side may be carried as the convention.

The next step separated sign from magnitude. The positive interior can read the
size of the residue, and that is real work, but it cannot tell the apparatus
which orientation to name as negative. Magnitude has no hidden arrow inside it.
To ask a magnitude for a sign is to ask the wrong instrument to answer the last
question. The chapter therefore refused a tempting shortcut. It would not let a
positive reading pretend to settle an orientation it is not built to see.

After that refusal, occurrence had to become a finite question. The apparatus
can count how often the two orientations appear. It can also specify a floor, a
discipline for close calls, so that a small excess is not mistaken for a
settled result. One orientation dominates only when it clears the other past
that floor. If neither side clears the floor, no anchor is produced. This is
not indecision dressed up as rigor. It is the price of letting a finite record
say only what it has earned the right to say.

S04 then made the anchor sharp. For a fixed tally and a fixed floor, both
opposite orientations cannot dominate at once. The rule has a one-place shape.
When an anchor exists, it gives the apparatus a tally, a floor, a dominant
orientation, and a margin. A later reader can inspect those four pieces without
sharing the apparatus's taste. The certificate is finite, and its uniqueness is
finite too: one record, one floor, at most one dominant side.

S05 named the consequence. The apparatus can anchor a convention by
measurement without deriving the sign. That sentence is the hinge on which the
chapter closes. Convention is still convention because the interior did not
force the sign. Measurement is still measurement because the tally, floor, and
margin give the convention a receipt. The chapter has neither dissolved the
sign into whim nor promoted a tally into an origin story. It has found the
middle position the finite apparatus can occupy with a clear conscience.

Now the needle can be named carefully. The needle is the decision edge at which
the apparatus seats a sign convention: which side of the residue it will carry
as negative. But this does not mean the first mark secretly contained charge
sign, or that every earlier use of the needle was already this final name in
disguise. Earlier chapters used the needle more broadly for the classical edge
the apparatus could expose but not compute away. Here, at the end of the
numbered argument, that broad edge receives its sharpest instance. The final
needle is the orientation convention, seated by a finite tally when the record
licenses it.

That distinction protects the whole arc from false hindsight. The book did not
begin with a hidden answer and spend seventeen chapters discovering what it had
already assumed. It began with finite marks, receipts, residues, energies,
invariants, boundaries, gauges, readings, and completion doors. At each stage it
asked what the apparatus could construct and what remained open. The sign
question is not retrofitted into all of those earlier decisions. It is the last
place where the same discipline becomes unmistakable: construct what can be
constructed, then name the convention only as far as the record permits.

The bridge therefore closes, rather than expands, the chapter's claim. The
apparatus has not explained why the measured occurrences leaned one way. It has
not found a hidden sign inside positive magnitude. It has not replaced a future
account of occurrence with a bookkeeping rule. What it has done is narrower and
stronger: it has shown how a finite record can license a carried sign
convention. Tally, floor, margin, record; these are the credentials. Without
them the sign remains unseated. With them the apparatus may carry the
orientation honestly.

The coda can now ask which pole the needle names. That question belongs to an
image rather than to a new proof. The argument itself is complete: opposite
orientations have been separated, magnitude has been refused as a sign oracle,
dominance has been made finite, uniqueness has been proved for a fixed record,
and convention has been anchored without being derived. The last bridge leaves
the reader there, at the seated needle, with the finite apparatus neither
silent nor grandiose. It has counted what it can count, carried what that count
licenses, and stopped at the boundary where explanation would have to come from
somewhere beyond this apparatus.

\section*{Coda: Which Pole the Needle Names}

A compass closes the chapter because it is a small instrument with a severe
honesty. It does not know the earth. It does not contain the field it answers.
It does not explain why the field runs as it does. It holds a needle, lets that
needle turn, waits for the motion to settle, and then names an orientation. A
hand can carry the compass because the compass does not pretend to carry the
whole cause of its reading. It carries only the disciplined surface where a
direction becomes visible.

The first lesson is that the needle reads direction, not strength. A magnetic
field has intensity as well as direction, but the compass needle is not a meter
for that intensity. It can swing sharply in a strong field and weakly in a weak
one, but its ordinary promise is not to quantify the field's size. It tells
which way. Another instrument may be needed for how much. Another account may
be needed for why. The needle's gift is orientation.

That is why the compass belongs here. Chapter 17 has just separated magnitude
from sign. The positive interior can report size. It can say that the residue
has unit magnitude. It can support comparison, detection, and finite
accounting. But it cannot say which side of the residue should be named
negative. The compass makes the distinction tactile. A hand may feel the case,
see the dial, and watch the needle settle; still, the needle is not telling the
strength of the field. It is showing the direction in which the reading has
come to rest.

The second lesson is that a pole name is a convention. The painted end of the
needle is called north-seeking because that is how the instrument has been
settled and taught. The word "north" is not engraved inside the bare physics of
the needle as a metaphysical necessity. The two ends of the needle are opposite
orientations of one object. The name attaches to one end because the instrument,
used in the world, gives a stable way to carry that end as the named one. The
convention is not whim, because the settled orientation supports it. It is not
derivation, because the name is still a name.

This is the coda's version of the chapter's central sentence: convention is
gauge, and tally is measurement. The compass has a symmetry of ends. If a child
painted the other end, the material needle would still have two opposite
directions. The field would not be rewritten by the paint. Yet once the
instrument is calibrated and its use is shared, the painted end gives a public
convention. A traveler does not need to decide north anew at every step. The
needle has anchored a label well enough to guide the hand.

But the anchoring depends on settlement. A compass held beside iron, shaken in a
moving vehicle, or placed where the field is weak may quiver without giving a
trustworthy direction. The needle moves, but it does not yet name. It swings
across the mark, hesitates, returns, and refuses to become a record. A careful
reader does not promote that trembling into a direction merely because one side
was ahead for an instant. The instrument has not cleared its own tremor.

The noise floor of the chapter is that tremor made explicit. A tally that gives
five readings on one side and four on the other may look like a preference if
one is impatient. The apparatus is not allowed to be impatient. It asks whether
the excess clears the floor. If it does not, the coda's compass is still
shaking. There is a motion, but not a settled direction; a count, but not a
licensed convention. The refusal is not failure. It is what keeps a finite
instrument from mistaking its own jitter for the world's answer.

That refusal is also what makes the eventual naming trustworthy. A compass that
admits its trembling can later be believed when it stops trembling. The pause is
part of the instrument's character. It teaches the hand not to treat every
motion as a reading. In the same way, the finite apparatus earns its convention
by having a visible no-anchor case. If the record is too close, it says so. If
the record is decisive, the later yes has weight because the apparatus has
already shown that it knows how to say no.

When the needle does settle, the convention becomes usable. The compass does
not need to repeat the whole history of magnetism each time it points. It gives
a present orientation under present conditions. The traveler may carry the
reading because the instrument has produced a stable enough mark. Likewise, the
apparatus may carry a sign convention when the record gives a tally, a floor, a
dominant orientation, and a margin. The convention remains a convention, but it
has a receipt. The needle points; the tally certifies.

A usable convention is not a universal explanation. A map may place north at
the top, a compass may mark which way north lies, and a traveler may use both
without confusing either with the earth itself. The usefulness of the convention
comes from being carried consistently after it has been anchored. It lets later
sentences, later diagrams, and later motions speak the same way. The chapter's
sign convention has that same modest power. Once the apparatus has a certified
orientation, it can speak with a settled sign. The settlement helps use; it does
not convert use into origin.

The two ends of the compass needle also keep the mirror honest. There is not a
north needle and a separate south needle competing inside the case. There is one
needle with two ends. Read from one end it gives one orientation; read from the
other it gives the opposite. The convention chooses which end to paint and name,
but the opposite end remains present. It is not destroyed by the name, and it is
not rendered unreal by losing the convention. It is the other side of the same
oriented object.

So too with the electron and positron readings in the chapter. They are not two
unrelated signs stored in different compartments of the apparatus. They are one
unit residue read from opposite sides. Calling one side negative does not erase
the other side. It seats a convention for carrying the sign. If the tally
anchors the electron orientation, the positron orientation remains the mirror
reading. The compass needle lets the reader feel this: one object, two ends,
one convention for naming which end guides the hand.

The compass also leaves the cause of the field outside its authority. It may
point steadily north on a table, but that steadiness is not an explanation of
the field. The needle does not know why the earth has the magnetic pattern it
has. It does not know the history of the ground beneath it, the material
condition of the needle, or the larger physical account that would explain why
the field arrives at the instrument with this direction rather than another.
The compass receives the field and reads orientation.

This outside cause is not an embarrassment to the compass. It is the reason the
compass can be honest. If the instrument claimed to contain the cause of the
field, it would be pretending to be much larger than it is. Its dignity is
smaller and cleaner: it stands at the meeting of needle and field and reports
the orientation available there. The apparatus's sign anchor has the same
dignity. It stands at the meeting of finite record and convention. It does not
swallow the deeper question. It gives the deeper question a boundary.

This restraint is exactly the restraint Chapter 17 has earned. The apparatus
can receive the world's occurrences of the two orientations. It can count them.
It can refuse close calls. It can certify a dominant orientation when the margin
clears the floor. It cannot explain why the occurrences leaned that way. The
lean is the field arriving at the compass. The tally is the needle settling
under that arrival. The convention is the painted end the instrument is then
licensed to carry.

The image should not be made heavier than that. This coda is not a theory of
magnetism. It is a way of holding the chapter's structure in the hand. A
compass has an interior arrangement, a needle, a dial, and a habit of use, but
the coda needs only its discipline: direction rather than strength, naming
rather than derivation, settlement past tremor, two ends of one object, and a
field whose cause remains outside the small instrument's authority. The image
works because it is limited.

There is a final discipline in the hand that holds the compass. The hand does
not make the needle point by wanting a direction. It steadies the case, keeps it
away from obvious disturbance, waits, and then reads. This is the posture the
book has asked of the reader from the first mark onward. Do not force the
apparatus to answer too early. Do not rename a wobble as a theorem. Do not ask a
small instrument to replace the world. Hold it well enough for its actual
reading to appear, and then carry exactly that reading.

With that limit, the chapter's last picture can be plain. The finite apparatus
has spent the book refusing to claim what it has not measured. Here it refuses
one final temptation. It does not say that the sign was hidden in magnitude. It
does not say that an interior calculation chose the electron. It does not say
that the count explains its own distribution. It says: the needle has settled
past the tremor; this end may be named; here is the record by which the name is
carried.

The compass in the hand is therefore the final image of a seated convention. It
is neither a guess nor an oracle. It is a finite instrument letting an
orientation become stable enough to name. The hand still belongs to a reader,
the field still belongs to the world, and the cause still lies deeper than the
needle. But the needle has done its work. It has turned, settled, and made one
pole nameable. That is all the apparatus ever promised at the sign: not the
origin of the world's lean, not a theorem forcing the name, but a measured
permission to carry one orientation as the convention.
